\label{sec:risk}
In this section, we describe the risks identified during this project.
In the previous section (\Cref{sec:author_evaluation}), we mainly identified issues related to the papers released in this study. In this section, we present various risks encountered during the development process.
Sharing such risks is essential to prevent overreliance on these systems and to promote a deeper understanding of AI Scientists within the research community.



\subsection{Idea Generation}
\begin{tcolorbox}[colback=gray!10, colframe=gray!40!black, title=Idea Risk1:,breakable]
     Identifying a successful idea is highly computationally expensive.
\end{tcolorbox}
The ideas generated by AI do not always work, which holds true for human scientists.
In our case, we aimed to generate one successful idea for each baseline paper.
To this end, we generated approximately ten ideas and evaluated them. Some were filtered out through human review, while others did not outperform the baseline. Finally, only one idea proved to be successful.

From this perspective, more extensive validation was conducted in some recent works \citep{liu2025alphago, weng2025deepscientist}.
For example, the concurrent work DeepScientist~\citep{weng2025deepscientist} performed a comprehensive large-scale study.
They report that, out of approximately 5,000 unique scientific ideas generated, only 21 ultimately led to genuine scientific innovations. Our experiments require less time because our limitation analysis is effective, and our goal is modest, aiming to find one successful idea rather than exploring more successful ideas.

Validating large-scale ideas is highly computationally expensive and often infeasible for many academic laboratories.
Future research will therefore focus on developing more efficient idea-pruning mechanisms, an efficient tree search algorithm for experiments, or incorporating human feedback.

\subsection{Experiment}
\begin{tcolorbox}[colback=gray!10, colframe=gray!40!black, title=Experiment Risk1 :,breakable]
     Lacking domain expertise, the coding agent sometimes produces code that leads to incorrect implementations and false performance gains.
\end{tcolorbox}
Because the coding agent is unaware of domain-specific conventions, it often improves performance in undesirable or invalid ways.
This issue frequently appeared in the experiments on GL-MCM~\citep{miyai2025gl}, which we describe in detail below.

\textbf{Background.}
GL-MCM is a task of zero-shot out-of-distribution (OOD) detection~\citep{miyai2024generalized}.
OOD detection aims to distinguish between samples belonging to a predefined in-distribution (ID) class set (e.g., the 1000 classes of ImageNet) and those belonging to classes with different semantics~\citep{yang2024generalized}.
In the GL-MCM setting, the model uses CLIP~\citep{radford2021learning} and is required to discriminate between ID and OOD data without any training, given only the ID class names.

As a convention in this research area, the source code is typically written as shown in \Cref{code:glmcm}.
Specifically, the ID and OOD dataloaders are defined separately.
A batch is first sampled from the ID dataloader to obtain an OOD score, followed by another batch from the OOD dataloader to compute its OOD score.
Finally, the OOD scores and the corresponding ID/OOD labels are used to compute the AUROC.

\begin{algorithm}[t]
\caption{GL-MCM implementation (Python-like pseudocode).}
\begin{algorithmic}[1]
\Require ID dataloader $\mathcal{D}_{ID}$, OOD dataloader $\mathcal{D}_{OOD}$, method $f$, scoring function $S(\cdot)$
\Ensure AUROC value
\State scores = []
\State labels = []
\For{batch in $\mathcal{D}_{ID}$}
  \State ood\_score = $S(f(\text{batch}))$
  \State scores.append(ood\_score)
  \State labels.append(0) \Comment{0 = ID sample}
\EndFor
\For{batch in $\mathcal{D}_{OOD}$}
  \State ood\_score = $S(f(\text{batch}))$
  \State scores.append(ood\_score)
  \State labels.append(1) \Comment{1 = OOD sample}
\EndFor
\State auroc = AUROC(scores, labels)
\State \Return auroc
\end{algorithmic}
\label{code:glmcm}
\end{algorithm}


\textbf{Mistake by Jr. AI Scientist.}
Our Jr. AI Scientist wrote code that applied bach-level normalization and statistical operations within the method $f$ for each batch.
However, as shown in \Cref{code:glmcm}, each batch contains only ID or OOD samples, not both. As a result, the batch-level statistics are biased toward either the ID or OOD distribution.
Human experts can immediately recognize that normalization should not be performed on a per-batch basis.
Nevertheless, during numerous attempts to improve performance, the Jr. AI Scientist often arrived at such invalid solutions.
We believe this issue will persist even as the performance of coding agents continues to improve.
This observation highlights the importance of human researchers possessing sufficient domain expertise to verify whether the observed performance improvements are indeed valid.

\subsection{Writing}

\begin{tcolorbox}[colback=gray!10, colframe=gray!40!black, title=Writing Risk1 :,breakable]
    When feedback is provided, fabrication of experimental results can easily occur.
\end{tcolorbox}
We found that feedback can sometimes become a major source of fabrication.
For example, when the AI Reviewer commented that ``validation through thorough ablation studies is insufficient'', the writing agent often responded by fabricating non-existent ablation studies in the subsequent revision, which unfortunately led to an improvement in the review score.
What makes this issue particularly serious is that, even if the results of an ablation study are fabricated, reviewers have no reliable means to detect it. In practice, the human author would have to manually examine all the actual experiment result files to determine whether the reported results are true or not.

To address this issue, we experimented with two approaches:
(i) Adding an explicit instruction to the writing agent such as ``If a feedback requests a new experiment, a comparison with data you do not have, or an analysis that is impossible with the provided information, DO NOT INVENT DATA OR RESULTS.'' to explicitly prohibit fabrication or falsification. (ii) Providing the writing agent with experimental results in a structured summary format that was both easy to parse and contained detailed descriptions of each setting and its corresponding outcomes. 
The second approach proved particularly important. Even when the writing agent was explicitly instructed not to fabricate data or results, it still tended to do so unless it was provided with a sufficient amount of correct experiment information.
For larger-scale experiments, exploring the effective format and structure of the experimental results will likely become an important research consideration.

Despite these improvements, hallucinations still occur, as shown in \Cref{sec:author_evaluation}. Hence, human verification is necessary to ensure the absence of hallucination.


\begin{tcolorbox}[colback=gray!10, colframe=gray!40!black, title=Writing Risk2 :,breakable]
    Making appropriate citations in the right context remains challenging.
\end{tcolorbox}

In our system, making appropriate citations in the right context still remains challenging. Through several design improvements, (i) we have prevented the agent from citing non-existent papers, and (2) it can correctly handle citations to papers included in the baseline.
However, issues remain with newly added BibTeX entries.
The agent sometimes cites these papers in irrelevant contexts.
This problem mainly arises from the current framework, in which the agent searches for related papers through the Semantic Scholar API, extracts their BibTeX entries and abstracts, summarizes them, and refers to these summaries when writing the manuscript.
Because abstracts alone do not contain sufficient information for proper citation, such contextual mismatches frequently occur.
Therefore, enabling an AI system to make appropriate citations likely requires a deeper, human-level understanding of the referenced papers, which remains a highly challenging problem.

\begin{tcolorbox}[colback=gray!10, colframe=gray!40!black, title=Writing Risk3 :,breakable]
    The result interpretation is unreliable.
\end{tcolorbox}
We found that the writing agent sometimes makes unreliable or unfounded interpretations of the results.
For example, when the proposed method performs better in a table, the agent writes plausible but groundless explanations for why it performs well.
Similarly, when referring to figures, it tends to exaggerate the effectiveness of the method beyond what can actually be seen.
This shows that accurately interpreting experimental results is still a difficult task for our AI Scientist system.

\begin{tcolorbox}[colback=gray!10, colframe=gray!40!black, title=Writing Risk4 :,breakable]
    A mechanism is needed to prevent the agent from generating non-existent citations.
\end{tcolorbox}
During the reflection stage, we observed that the agent occasionally modified the BibTeX file on its own—for example, by introducing incorrect author information or adding entries for papers that do not actually exist.
To address this issue, we adopted an agentic framework in which, whenever a citation is required during feedback-based revision, any references to be revised or added are dynamically retrieved through the Semantic Scholar API.
In addition, since the writing agent sometimes automatically generated a new .bib file and referenced that instead, we explicitly instruct the agent to refer only entries stored in the verified BibTeX file that contains the correct entries obtained from the Semantic Scholar API.



\subsection{Review}
\begin{tcolorbox}[colback=gray!10, colframe=gray!40!black, title=Review Risk1 :,breakable]
     Current AI reviewers cannot detect discrepancies between the actual experimental results and the written descriptions.
\end{tcolorbox}

Current AI Reviewers primarily evaluate the written content of papers and lack any mechanism to detect discrepancies between the text and the actual experimental results.
For instance, even if all the reported ablation studies were fabricated, there is no way for the reviewer to identify such inconsistencies. A similar observation was also reported in \citep{jiang2025badscientist}.
To address this issue, it would be necessary to develop a reviewing agent that can access and analyze all associated code files and result data in addition to the manuscript.
Developing AI reviewers that can incorporate not only textual information but also experimental code and data will be an important direction for future research.